\documentclass[11pt]{article}
% Page layout and typography
\usepackage[margin=1in]{geometry}
\usepackage{microtype}
\usepackage{parskip}

% Mathematics
\usepackage{amsmath,amssymb,amsthm,mathtools}
\usepackage{bm}

% Figures and links
\usepackage{graphicx}
\graphicspath{{figures/}}
\usepackage{xcolor}
\usepackage[hidelinks]{hyperref}

% Lists
\usepackage{enumitem}
\setlist{nosep}

% Theorem environments
\theoremstyle{definition}
\newtheorem{definition}{Definition}[section]
\newtheorem{example}[definition]{Example}
\newtheorem{remark}[definition]{Remark}

\theoremstyle{plain}
\newtheorem{theorem}[definition]{Theorem}
\newtheorem{lemma}[definition]{Lemma}
\newtheorem{proposition}[definition]{Proposition}
\newtheorem{corollary}[definition]{Corollary}

% Common number systems
\newcommand{\R}{\mathbb{R}}
\newcommand{\C}{\mathbb{C}}
\newcommand{\Q}{\mathbb{Q}}
\newcommand{\Z}{\mathbb{Z}}
\newcommand{\N}{\mathbb{N}}
\newcommand{\D}{\mathbb{D}}
\newcommand{\T}{\mathbb{T}}
\newcommand{\RS}{\mathbb{S}^2}
\newcommand{\per}{\text{per}}

% Common operators
\DeclareMathOperator{\im}{im}
\DeclareMathOperator{\rank}{rank}
\DeclareMathOperator{\Span}{span}
\DeclareMathOperator{\id}{id}


\title{Algebraic Numbers and Functions Course Notes}
\author{Caelen Gelineau}
\date{\today}

\begin{document}

\maketitle

\section{Galois review}
\begin{definition}
    Seperable means equivalently
    \begin{enumerate}
        \item the number of distinct $k$-embeddings $L \hookrightarrow K$ is $[L:K]$
        \item For $\alpha \in L$, the minimal polynomial of $\alpha$ over $K$ is "seperable": having no repeated 
        roots in $\overline{K}$, i.e. gcd($f(x), f'(x)$) $ = 1$.
    \end{enumerate}
    
\end{definition}
\begin{definition}
    Normal means equivalently
    \begin{enumerate}
        \item every $k$-embedding $L \hookrightarrow \overline{K}$ has the same image
        \item $L$ is a splitting field (of a family of polynomials).
        \item Given an irreducible polynomial $f(x) \in K[x]$, either has no roots in $L$ or all of its roots are in $L$.
    \end{enumerate}
    
\end{definition}

\begin{definition}
    $K \subseteq L$ is Galois if it is separable and normal.
\end{definition}

\begin{definition}
    \[
        \text{Gal}(L/K) := \{ \sigma: L \rightarrow L \text{ automorphisms s.t. } \sigma|_k = id_k\}
    \]
    (fixing base $k$-embedding $L \hookrightarrow \overline{K})$, Gal is identified with k-embeddings $L \hookrightarrow \overline{K}$.
\end{definition}
If $L$ is a finite Galois group over $K$, the Galois correspondence:
subgroups
\[
    \{H \le G\}
\]
intermediate fields:
\[
    \{K \le F \le L\}
\]
$H \mapsto L^H := \{\alpha \in L | \sigma(\alpha) = \alpha \text{ for } \sigma \in H\} $ and
\[
    F \mapsto \text{Gal}(L/F)
\]
$H \trianglelefteq G$ iff $F$ is Galois / $K$.\\
$|\text{Gal}(L/K)| = [L:K]$.



\section{PRime ideals}
$P$ is a prime ideal. We have $P \subset R \subset K$. If we Galois extension $K$ to $L$, we have $Q \subset S \subset L$.
\begin{definition}
    $P \subset Q$ means "$Q$ is a prime over $P$".
\end{definition}
\begin{theorem}
    Theorem 19 of marcus.
    \begin{enumerate}
        \item $P \subset Q \implies Q \cap R = P$.
    \end{enumerate}
    
\end{theorem}
\begin{theorem}
    Theorem 20 of marcus.\\
    Every $P$ has exactly $\ge 1$ $Q$ and every $Q$ has exactly 1 $P$.
\end{theorem}









\end{document}